\documentclass{article}
\usepackage[margin=1in]{geometry}
\usepackage[T1]{fontenc}
\usepackage{lmodern}
\usepackage{microtype}
\usepackage{graphicx} % Required for inserting images
\usepackage{amsmath,amssymb,amsthm}
\usepackage{mathtools} % \xRightarrow (labelled arrows)
% for theorem type-format, limits double subscript, etc.
\usepackage{enumitem} % compact / indented lists
\usepackage{float}    % [H] = put the figure exactly here, never float it away
\usepackage{needspace} % keep a "Step" heading on the same page as what follows it
\usepackage{titlesec} % control the space around section headings

% more air around section / subsection headings
\titlespacing*{\section}{0pt}{3.5em plus 1em minus .5em}{1.2em plus .3em}
\titlespacing*{\subsection}{0pt}{2.5em plus .8em minus .4em}{1em plus .2em}

\usepackage{tikz} % for drawings...
\usetikzlibrary{tikzmark}
\usetikzlibrary{calc}
\usetikzlibrary{decorations.markings} % going x% along bezier curve
\usetikzlibrary{patterns} % for hatched shading
\usetikzlibrary{decorations.pathreplacing} % curly braces in diagrams

\usepackage{mathrsfs} % \mathscr{L} % fancy script L
\usepackage[hidelinks]{hyperref} % clickable \ref's, no coloured boxes

\usetikzlibrary{arrows.meta} %for stealth arrow tips

\title{v5 construction of R}
\author{Lex S}
\date{October 2026}


% how do pt vs nothing-units compare?? --- 5 pt=0.176 cm --- ish
% basically pt means in units of letters size
% and nothing units is basically in cm

%\    1 word space    a\ b
%\quad    1 em    a\quad b
%\qquad    2 em    a\qquad b

% spacing: a little more air everywhere
\setlength{\jot}{6pt}                    % space between the rows of an align (default 3pt)
\setlength{\parskip}{2pt plus 1pt}       % a hair of space between paragraphs
\makeatletter
% same look as amsthm's "definition" style, but 12pt above and below instead of 8pt
\newtheoremstyle{spaced}{12pt plus 2pt minus 2pt}{12pt plus 2pt minus 2pt}{}{}{\bfseries}{.}{5pt plus 1pt minus 1pt}{}
% amsthm's proof environment with 10pt (instead of 6pt) above and below
\renewenvironment{proof}[1][\proofname]{\par
  \pushQED{\qed}%
  \normalfont \topsep 10\p@\@plus4\p@\relax
  \trivlist
  \item[\hskip\labelsep\itshape #1\@addpunct{.}]\ignorespaces
}{\popQED\endtrivlist\@endpefalse}
\makeatother

\theoremstyle{spaced}
\newtheorem{definition}{DEFINITION}[section]
\newtheorem{theorem}[definition]{THEOREM}
\newtheorem{lemma}[definition]{Lemma}
\newtheorem{corollary}[definition]{Corollary}
\newtheorem{remark}[definition]{Remark}
\newtheorem{example}[definition]{Example}

% shorthands
\newcommand{\Q}{\mathbb{Q}}
\newcommand{\N}{\mathbb{N}}
\newcommand{\Z}{\mathbb{Z}}
\newcommand{\R}{\mathbb{R}}
\newcommand{\Qb}{\bar{\mathbb{Q}}}
\newcommand{\db}{\bar{d}}
\newcommand{\Vb}{\bar{V}}

% proof-chain notation
\newcommand{\imp}{\;\Rightarrow\;}                                   % implication
\newcommand{\by}[1]{\;\xRightarrow{\,\text{\scriptsize #1}\,}\;}     % implication, with the reason on the arrow
\newcommand{\byiff}[1]{\;\xLeftrightarrow{\,\text{\scriptsize #1}\,}\;}
\newcommand{\ev}{\ (n\gg0)}                                          % "for all sufficiently large n"
\newcommand{\step}[2]{\par\Needspace{8\baselineskip}\vspace{1.5em}\noindent\textbf{Step #1 $\mid$} #2\par}
\newcommand{\substep}[2]{\par\vspace{1em}\noindent\hspace{1.5em}\textbf{#1 $\mid$} #2}

